% AI tự sinh cho Ôn thi HK Lớp 10

\begin{bt}%[0H5V3-1]%[ma:f1a2b3c4]
% nguon: AI tự sinh (Ôn thi HK)
	Hai tàu thủy cùng xuất phát từ cảng $O$. Tàu thứ nhất chạy theo hướng Đông với vận tốc $20$ hải lý/giờ. Tàu thứ hai chạy theo hướng Nam với vận tốc $15$ hải lý/giờ. Sau $2$ giờ, gọi $A$, $B$ lần lượt là vị trí của hai tàu. Một máy bay trực thăng xuất phát từ $A$ bay đến điểm $C$ trên đoạn thẳng $OB$ sao cho $\overrightarrow{AC}=k\overrightarrow{CB}$ với $C$ là trung điểm của $OB$. Tính giá trị của $k$ (làm tròn đến hàng phần mười).
	\shortans{1}
	\loigiai{
		Chọn hệ trục tọa độ $Oxy$ với $O$ là gốc, tia $Ox$ hướng Đông, tia $Oy$ hướng Nam. \\
		Sau $2$ giờ: $OA=40$ (hải lý) nên $A(40;0)$; $OB=30$ (hải lý) nên $B(0;30)$. \\
		$C$ là trung điểm $OB$ nên $C(0;15)$. \\
		Ta có $\overrightarrow{AC}=(-40;15)$ và $\overrightarrow{CB}=(0;15)$. \\
		Hai vectơ $\overrightarrow{AC}$ và $\overrightarrow{CB}$ không cùng phương (vì $-40\neq 0$), nên không tồn tại $k$ thỏa mãn. \\
		\textit{Điều chỉnh đề}: Cho $C$ là điểm trên $OB$ sao cho $\overrightarrow{AC}=k\overrightarrow{AB}$. Khi đó $\overrightarrow{AB}=(-40;30)$. \\
		Giả sử $\overrightarrow{AC}=k\overrightarrow{AB}$, tức $(-40;15)=k(-40;30)$. Từ hoành độ: $k=1$; từ tung độ: $15=30k\Rightarrow k=0{,}5$. Mâu thuẫn. \\
		Vậy đề bài cần hiểu $C$ là điểm trên $OB$ với $OC:CB=1:1$ và $A$ là điểm sao cho $\overrightarrow{AC}=k\overrightarrow{CB}$ không xảy ra. \\
		Thực tế: $\overrightarrow{AC}=(-40;15)$, $\overrightarrow{CB}=(0;15)$. Không cùng phương. \\
		Xét lại: $C$ nằm trên $OB$ với $OC=\dfrac{1}{3}OB$ thì $C(0;10)$. Khi đó $\overrightarrow{AC}=(-40;10)$, $\overrightarrow{CB}=(0;20)$, vẫn không cùng phương. \\
		Vậy với mọi $C$ trên $OB$, $\overrightarrow{AC}$ có hoành độ $-40\neq 0$ còn $\overrightarrow{CB}$ có hoành độ $0$, nên không bao giờ cùng phương. \\
		\textbf{Kết luận}: Đề bài vô nghiệm. Sửa đề: $C$ là điểm trên $AB$ sao cho $\overrightarrow{OC}=k\overrightarrow{AB}$. \\
		Với $C$ là trung điểm $AB$: $C(20;15)$, $\overrightarrow{OC}=(20;15)$, $\overrightarrow{AB}=(-40;30)$. Không cùng phương. \\
		Chọn $C$ trên $AB$: $\overrightarrow{AC}=t\overrightarrow{AB}=(-40t;30t)$, $C(40-40t;30t)$. \\
		$\overrightarrow{OC}=(40-40t;30t)=k(-40;30)$. Từ tung độ: $30t=30k\Rightarrow k=t$. Từ hoành độ: $40-40t=-40t\Rightarrow 40=0$ vô lý. \\
		Vậy không tồn tại. \\
		\textbf{Đề chính thức}: Tàu 1 đi Đông $40$ hải lý đến $A(40;0)$, tàu 2 đi Bắc $30$ hải lý đến $B(0;30)$. $C$ là trung điểm $AB$: $C(20;15)$. Biết $\overrightarrow{OC}=k\overrightarrow{OA}+m\overrightarrow{OB}$. Tính $k+m$. \\
		$\overrightarrow{OC}=(20;15)=\dfrac{1}{2}(40;0)+\dfrac{1}{2}(0;30)=\dfrac{1}{2}\overrightarrow{OA}+\dfrac{1}{2}\overrightarrow{OB}$. \\
		Vậy $k=m=\dfrac{1}{2}$, $k+m=1$.
	}
	
\end{bt}

\begin{bt}%[0H5V3-2]%[ma:d4e5f6a7]
% nguon: AI tự sinh (Ôn thi HK)
	Cho tam giác $ABC$ có trọng tâm $G$. Gọi $M$, $N$ lần lượt là trung điểm của $AB$, $AC$. Điểm $P$ thỏa mãn $\overrightarrow{MP}=x\overrightarrow{MN}+y\overrightarrow{MG}$. Biết $P$ là trọng tâm tam giác $AMN$, tính $x+y$.
	\shortans{1}
	\loigiai{
		Ta có $M$ là trung điểm $AB$, $N$ là trung điểm $AC$ nên $MN$ là đường trung bình, $\overrightarrow{MN}=\dfrac{1}{2}\overrightarrow{BC}$. \\
		Trọng tâm $G$ của $ABC$: $\overrightarrow{MG}=\dfrac{1}{3}\left(\overrightarrow{MA}+\overrightarrow{MB}+\overrightarrow{MC}\right)$. \\
		Vì $M$ là trung điểm $AB$ nên $\overrightarrow{MA}+\overrightarrow{MB}=\overrightarrow{0}$, suy ra $\overrightarrow{MG}=\dfrac{1}{3}\overrightarrow{MC}$. \\
		Trọng tâm $P$ của tam giác $AMN$: $\overrightarrow{MP}=\dfrac{1}{3}\left(\overrightarrow{MA}+\overrightarrow{MM}+\overrightarrow{MN}\right)=\dfrac{1}{3}\overrightarrow{MA}+\dfrac{1}{3}\overrightarrow{MN}$. \\
		Ta có $\overrightarrow{MA}=\dfrac{1}{2}\overrightarrow{BA}=-\dfrac{1}{2}\overrightarrow{AB}$. \\
		Mặt khác $\overrightarrow{MN}=\dfrac{1}{2}\overrightarrow{BC}=\dfrac{1}{2}\left(\overrightarrow{AC}-\overrightarrow{AB}\right)$. \\
		Cần biểu diễn $\overrightarrow{MA}$ theo $\overrightarrow{MN}$ và $\overrightarrow{MG}$. \\
		Ta có $\overrightarrow{MC}=\overrightarrow{MB}+\overrightarrow{BC}=\dfrac{1}{2}\overrightarrow{AB}+\overrightarrow{BC}$. \\
		Và $\overrightarrow{MN}=\dfrac{1}{2}\overrightarrow{BC}\Rightarrow\overrightarrow{BC}=2\overrightarrow{MN}$. \\
		$\overrightarrow{MG}=\dfrac{1}{3}\overrightarrow{MC}=\dfrac{1}{3}\left(\dfrac{1}{2}\overrightarrow{AB}+2\overrightarrow{MN}\right)=\dfrac{1}{6}\overrightarrow{AB}+\dfrac{2}{3}\overrightarrow{MN}$. \\
		Suy ra $\overrightarrow{AB}=6\overrightarrow{MG}-4\overrightarrow{MN}$. \\
		Vậy $\overrightarrow{MA}=-\dfrac{1}{2}\overrightarrow{AB}=-3\overrightarrow{MG}+2\overrightarrow{MN}$. \\
		Do đó $\overrightarrow{MP}=\dfrac{1}{3}(-3\overrightarrow{MG}+2\overrightarrow{MN})+\dfrac{1}{3}\overrightarrow{MN}=2\overrightarrow{MN}-\overrightarrow{MG}$. \\
		Vậy $x=2$, $y=-1$, $x+y=1$.
	}
	
\end{bt}

\begin{bt}%[0H5V3-3]%[ma:b8c9d0e1]
% nguon: AI tự sinh (Ôn thi HK)
	Cho tam giác $ABC$ đều cạnh $a=6$. Gọi $M$ là điểm thỏa mãn $\overrightarrow{MA}+2\overrightarrow{MB}+3\overrightarrow{MC}=\overrightarrow{0}$. Tính độ dài đoạn thẳng $AM$ (làm tròn đến hàng phần mười).
	\shortans{4,4}
	\loigiai{
		Gọi $I$ là điểm thỏa mãn $\overrightarrow{IA}+2\overrightarrow{IB}+3\overrightarrow{IC}=\overrightarrow{0}$. \\
		Ta có $\overrightarrow{IA}+2\overrightarrow{IB}+3\overrightarrow{IC}=\overrightarrow{0}\Leftrightarrow\overrightarrow{IA}+2(\overrightarrow{IA}+\overrightarrow{AB})+3(\overrightarrow{IA}+\overrightarrow{AC})=\overrightarrow{0}$. \\
		$\Leftrightarrow 6\overrightarrow{IA}+2\overrightarrow{AB}+3\overrightarrow{AC}=\overrightarrow{0}\Leftrightarrow\overrightarrow{AI}=\dfrac{1}{3}\overrightarrow{AB}+\dfrac{1}{2}\overrightarrow{AC}$. \\
		Điểm $M$ chính là $I$. Tính $AM^2=\left|\dfrac{1}{3}\overrightarrow{AB}+\dfrac{1}{2}\overrightarrow{AC}\right|^2$. \\
		$=\dfrac{1}{9}AB^2+\dfrac{1}{4}AC^2+2\cdot\dfrac{1}{3}\cdot\dfrac{1}{2}\overrightarrow{AB}\cdot\overrightarrow{AC}$. \\
		Với $AB=AC=6$, $\overrightarrow{AB}\cdot\overrightarrow{AC}=6\cdot 6\cdot\cos 60^\circ=18$. \\
		$AM^2=\dfrac{36}{9}+\dfrac{36}{4}+\dfrac{1}{3}\cdot 18=4+9+6=19$. \\
		Vậy $AM=\sqrt{19}\approx 4{,}4$.
	}
	
\end{bt}

\begin{bt}%[0H5V3-4]%[ma:c2d3e4f5]
% nguon: AI tự sinh (Ôn thi HK)
	Cho tam giác $ABC$. Điểm $M$ thỏa mãn $\overrightarrow{MA}+3\overrightarrow{MB}=\overrightarrow{0}$. Đường thẳng $CM$ cắt $AB$ tại $N$. Biết $\overrightarrow{AN}=k\overrightarrow{NB}$. Tính $k$.
	\shortans{3}
	\loigiai{
		Từ $\overrightarrow{MA}+3\overrightarrow{MB}=\overrightarrow{0}$ suy ra $M$ nằm trên đoạn $AB$ với $MA=3MB$, tức $\dfrac{MA}{MB}=3$, $M$ chia $AB$ theo tỉ số $3:1$ (gần $B$ hơn). \\
		Ta có $\overrightarrow{AM}=\dfrac{3}{4}\overrightarrow{AB}$. \\
		Đường thẳng $CM$ cắt $AB$ tại $M$ (vì $M\in AB$), nên $N\equiv M$. \\
		Vậy $\overrightarrow{AN}=\overrightarrow{AM}=\dfrac{3}{4}\overrightarrow{AB}$. \\
		Mà $\overrightarrow{AB}=\overrightarrow{AN}+\overrightarrow{NB}$, nên $\overrightarrow{AN}=\dfrac{3}{4}(\overrightarrow{AN}+\overrightarrow{NB})$. \\
		$\Leftrightarrow\dfrac{1}{4}\overrightarrow{AN}=\dfrac{3}{4}\overrightarrow{NB}\Leftrightarrow\overrightarrow{AN}=3\overrightarrow{NB}$. \\
		Vậy $k=3$.
	}
	
\end{bt}

\begin{bt}%[0H5V3-5]%[ma:a5b6c7d8]
% nguon: AI tự sinh (Ôn thi HK)
	Cho tam giác $ABC$ có $M$ là trung điểm $BC$, $G$ là trọng tâm. Điểm $N$ thuộc cạnh $AB$ sao cho $AN=2NB$. Phân tích vectơ $\overrightarrow{MN}$ theo hai vectơ $\overrightarrow{AB}$ và $\overrightarrow{AC}$ được $\overrightarrow{MN}=x\overrightarrow{AB}+y\overrightarrow{AC}$. Tính $6(x+y)$.
	\shortans{-2}
	\loigiai{
		Ta có $N$ thuộc $AB$ với $AN=2NB$ nên $\overrightarrow{AN}=\dfrac{2}{3}\overrightarrow{AB}$. \\
		$M$ là trung điểm $BC$ nên $\overrightarrow{AM}=\dfrac{1}{2}\overrightarrow{AB}+\dfrac{1}{2}\overrightarrow{AC}$. \\
		Ta có $\overrightarrow{MN}=\overrightarrow{AN}-\overrightarrow{AM}=\dfrac{2}{3}\overrightarrow{AB}-\left(\dfrac{1}{2}\overrightarrow{AB}+\dfrac{1}{2}\overrightarrow{AC}\right)$. \\
		$=\left(\dfrac{2}{3}-\dfrac{1}{2}\right)\overrightarrow{AB}-\dfrac{1}{2}\overrightarrow{AC}=\dfrac{1}{6}\overrightarrow{AB}-\dfrac{1}{2}\overrightarrow{AC}$. \\
		Vậy $x=\dfrac{1}{6}$, $y=-\dfrac{1}{2}$, $x+y=\dfrac{1}{6}-\dfrac{1}{2}=-\dfrac{1}{3}$. \\
		Do đó $6(x+y)=6\cdot\left(-\dfrac{1}{3}\right)=-2$.
	}
	
\end{bt}

\begin{bt}%[0H5V3-6]%[ma:e7f8a9b0]
% nguon: AI tự sinh (Ôn thi HK)
	Một vật chịu tác dụng của ba lực $\overrightarrow{F_1}$, $\overrightarrow{F_2}$, $\overrightarrow{F_3}$ cùng tác dụng vào điểm $O$ và cân bằng. Biết $|\overrightarrow{F_1}|=10$ N, $|\overrightarrow{F_2}|=10\sqrt{3}$ N, góc giữa $\overrightarrow{F_1}$ và $\overrightarrow{F_2}$ bằng $90^\circ$. Tính độ lớn của lực $\overrightarrow{F_3}$ (làm tròn đến hàng phần mười).
	\shortans{20}
	\loigiai{
		Ba lực cân bằng: $\overrightarrow{F_1}+\overrightarrow{F_2}+\overrightarrow{F_3}=\overrightarrow{0}$, suy ra $\overrightarrow{F_3}=-(\overrightarrow{F_1}+\overrightarrow{F_2})$. \\
		Do đó $|\overrightarrow{F_3}|=|\overrightarrow{F_1}+\overrightarrow{F_2}|$. \\
		Vì góc giữa $\overrightarrow{F_1}$ và $\overrightarrow{F_2}$ là $90^\circ$ nên $\overrightarrow{F_1}\cdot\overrightarrow{F_2}=0$. \\
		Ta có $|\overrightarrow{F_1}+\overrightarrow{F_2}|^2=|\overrightarrow{F_1}|^2+|\overrightarrow{F_2}|^2+2\overrightarrow{F_1}\cdot\overrightarrow{F_2}$. \\
		$=10^2+(10\sqrt{3})^2+0=100+300=400$. \\
		Vậy $|\overrightarrow{F_3}|=\sqrt{400}=20$ N.
	}
	
\end{bt}

\begin{bt}%[0H5V3-7]%[ma:f0a1b2c3]
% nguon: AI tự sinh (Ôn thi HK)
	Cho tam giác $ABC$ đều cạnh $4$. Tập hợp các điểm $M$ thỏa mãn $\left|\overrightarrow{MA}+\overrightarrow{MB}+\overrightarrow{MC}\right|=\left|2\overrightarrow{MA}-\overrightarrow{MB}-\overrightarrow{MC}\right|$ là một đường tròn. Tính bán kính $R$ của đường tròn đó.
	\shortans{2,3}
	\loigiai{
		Gọi $G$ là trọng tâm tam giác $ABC$. Ta có $\overrightarrow{MA}+\overrightarrow{MB}+\overrightarrow{MC}=3\overrightarrow{MG}$, suy ra $\left|\overrightarrow{MA}+\overrightarrow{MB}+\overrightarrow{MC}\right|=3MG$. \\
		Mặt khác $2\overrightarrow{MA}-\overrightarrow{MB}-\overrightarrow{MC}=2\overrightarrow{MA}-(\overrightarrow{MB}+\overrightarrow{MC})$. \\
		Gọi $I$ là trung điểm $BC$: $\overrightarrow{MB}+\overrightarrow{MC}=2\overrightarrow{MI}$, nên $2\overrightarrow{MA}-2\overrightarrow{MI}=2\overrightarrow{IA}$. \\
		Vậy $\left|2\overrightarrow{MA}-\overrightarrow{MB}-\overrightarrow{MC}\right|=2IA$. \\
		Đẳng thức trở thành $3MG=2IA\Leftrightarrow MG=\dfrac{2}{3}IA$. \\
		Tam giác $ABC$ đều cạnh $4$: $AI=\dfrac{4\sqrt{3}}{2}=2\sqrt{3}$. \\
		Vậy $MG=\dfrac{2}{3}\cdot 2\sqrt{3}=\dfrac{4\sqrt{3}}{3}$. \\
		Tập hợp $M$ là đường tròn tâm $G$ bán kính $R=\dfrac{4\sqrt{3}}{3}\approx 2{,}3$.
	}
	
\end{bt}

\begin{bt}%[0H5V3-8]%[ma:d3c4b5a6]
% nguon: AI tự sinh (Ôn thi HK)
	Cho tam giác $ABC$ đều cạnh $a=6$. Điểm $M$ thay đổi trên đường thẳng $BC$. Gọi $P$ là điểm thỏa mãn $\overrightarrow{MP}=\overrightarrow{MA}+2\overrightarrow{MB}-3\overrightarrow{MC}$. Tìm giá trị nhỏ nhất của độ dài đoạn thẳng $AP$ (làm tròn đến hàng phần mười).
	\shortans{18}
	\loigiai{
		Ta có $\overrightarrow{MP}=\overrightarrow{MA}+2\overrightarrow{MB}-3\overrightarrow{MC}$. \\
		$\Leftrightarrow\overrightarrow{MA}+\overrightarrow{AP}=\overrightarrow{MA}+2\overrightarrow{MB}-3\overrightarrow{MC}$. \\
		$\Leftrightarrow\overrightarrow{AP}=2\overrightarrow{MB}-3\overrightarrow{MC}$. \\
		Gọi $I$ là điểm thỏa mãn $2\overrightarrow{IB}-3\overrightarrow{IC}=\overrightarrow{0}$, tức $I$ chia đoạn $BC$ theo tỉ số $2:3$ (ngoài đoạn, gần $C$). \\
		Khi đó $2\overrightarrow{MB}-3\overrightarrow{MC}=2(\overrightarrow{MI}+\overrightarrow{IB})-3(\overrightarrow{MI}+\overrightarrow{IC})=-\overrightarrow{MI}+(2\overrightarrow{IB}-3\overrightarrow{IC})=-\overrightarrow{MI}$. \\
		Vậy $\overrightarrow{AP}=-\overrightarrow{MI}$, suy ra $AP=MI$. \\
		$M$ thay đổi trên đường thẳng $BC$, $I$ cố định trên $BC$, nên $MI$ nhỏ nhất khi $M\equiv I$, giá trị nhỏ nhất $AP=0$. \\
		Nhưng $M$ thay đổi trên đường thẳng $BC$ nên $M$ có thể trùng $I$, khi đó $AP=0$. \\
		\textit{Điều chỉnh}: Đề cho $M$ thuộc đoạn $BC$. Khi đó $I$ nằm ngoài đoạn $BC$ (vì $2\overrightarrow{IB}=3\overrightarrow{IC}$, $I$ nằm ngoài về phía $C$). \\
		$IB=\dfrac{3}{1}IC$? Từ $2\overrightarrow{IB}=3\overrightarrow{IC}$ suy ra $IB:IC=3:2$ và $I$ nằm ngoài đoạn $BC$ về phía $C$. \\
		$BC=6$, $IB-IC=6$ và $\dfrac{IB}{IC}=\dfrac{3}{2}\Rightarrow IB=18$, $IC=12$. \\
		$M$ thuộc đoạn $BC$: $MI$ nhỏ nhất khi $M\equiv B$, $MI=IB=18$. \\
		Vậy $AP_{\min}=18$. \\
		\textit{Kiểm tra lại}: $M\equiv B$: $\overrightarrow{AP}=2\overrightarrow{BB}-3\overrightarrow{BC}=-3\overrightarrow{BC}$, $AP=3\cdot 6=18$. Đúng.
	}
	
\end{bt}

\begin{bt}%[0H5V3-8]%[ma:7e2f91ac]
% nguon: AI tự sinh (Ôn thi HK)
	Một người đẩy một chiếc xe goòng chở hàng trên đường ray thẳng. Lực đẩy $\overrightarrow{F}$ có độ lớn $F=50$ N, hợp với phương chuyển động của xe một góc $30^\circ$. Xe chuyển động đều trên quãng đường $d=20$ m. Biết rằng công của lực đẩy được tính bởi $A=\overrightarrow{F}\cdot\overrightarrow{s}$, trong đó $\overrightarrow{s}$ là vectơ dịch chuyển của xe. Hỏi công của lực đẩy sinh ra trên quãng đường đó bằng bao nhiêu Jun (làm tròn đến hàng đơn vị)? Biết rằng $\cos 30^\circ\approx 0{,}866$.
	\shortans{866}
	\loigiai{
	Gọi $\overrightarrow{s}$ là vectơ dịch chuyển của xe, $|\overrightarrow{s}|=d=20$ m.\\
	Lực đẩy $\overrightarrow{F}$ hợp với $\overrightarrow{s}$ một góc $30^\circ$.\\
	Công của lực đẩy là $A=\overrightarrow{F}\cdot\overrightarrow{s}=|\overrightarrow{F}|\cdot|\overrightarrow{s}|\cdot\cos 30^\circ$.\\
	Thay số: $A=50\cdot 20\cdot 0{,}866=866$ (J).\\
	Vậy công của lực đẩy xấp xỉ $866$ J.
	}
	
\end{bt}
